\section{Setting and Controlled Protocol}
\label{sec:setting}
% A/B round 1: S1-S2 clearly better than LongMemEval 3.1-3.2; S3-S5 comparable to Lost in the Middle 2.1.
% Edits applied: fixed ingestion model, splits with seeds, call-count honesty, judge caveat, prompt numbers, decoding.

% S1: task
\paragraph{Task.}
A conversation is a sequence of dated sessions between a user and an assistant. A memory system reads the sessions in
order and maintains a store, calling an ingestion LLM and an embedding model as it needs. After the last session it
receives a question $q$ and the current date, builds a context $C(q)$ from its store, and a \emph{reader} model answers
$q$ from $C(q)$ under an answer prompt $P$. A judge model grades the answer. The score of a system therefore depends on
the memory through $C(q)$, but also on the reader, on $P$, on the length of $C(q)$, and on the judge. Our protocol
fixes the reader, $P$, and the judge, fixes the ingestion LLM and the embedding model across systems, and reports the
length of $C(q)$ for every system, so that two rows in a table differ in the memory, or in a single named factor
that we vary deliberately.

% S2: data
\paragraph{Data.}
BEAM~\citep{tavakoli2025beam} has 2,000 questions over 100 synthetic conversations in four length tiers. Every
session carries a date, and each conversation contains two questions for each of ten abilities: abstention,
contradiction resolution, event ordering, information extraction, instruction following, knowledge update,
multi-session reasoning, preference following, summarization, and temporal reasoning. We use the 100K tier
(20 conversations) and the 1M tier, where a conversation spans about 1,700 turns and 530,000 words, several times
more than the context we give any reader. We fixed the splits before running any experiment. On BEAM-100K, we
shuffle the 20 conversations with seed 0 and use the first ten (200 questions) for development and the other ten
(200 questions) for testing. On BEAM-1M, the first five conversations in the released order (100 questions) are
for development; of the remaining 30, we fixed the next ten (200 questions) for testing before running any test, because building every memory for all 30 would exceed our budget.
LongMemEval-S~\citep{wu2025longmemeval} gives each of its 500 questions its own history of about 115K tokens. We
stratify questions by type, treating abstention variants as their own stratum, and draw with seed 0 a development set
of \prelim{62} questions and a disjoint test set of 151; the remaining questions are unused. Every result before
\cref{sec:test} comes from development sets, and baselines are configured on the same development sets.

% S3: readers, judge, prompts
\paragraph{Readers, judge, and prompts.}
Our main reader is gpt-6-luna. To test whether conclusions depend on the reader, we repeat the main comparison with
DeepSeek-V4.1-Flash (weights of 2026-10-17, served with vLLM on eight H20 GPUs), with gpt-6.1-sol, and, for one run
because of its cost, with gpt-6-astra. Within a
comparison, every system uses the same reader, and every LLM call it makes at read time (query planning, conflict
checks, compilation) uses that reader too. Memory stores are built once with gpt-6-luna and shared, and Hindsight's
recalled contexts are computed once and reused; swapping the reader therefore tests the read side only. Answers and
compilation use medium reasoning effort; conflict checks and grading use low effort (thinking disabled for
DeepSeek). All answers are graded by gpt-6-luna with each benchmark's official prompt: on BEAM, a rubric prompt that
scores each rubric item 0, 0.5, or 1, averaged per question; on LongMemEval, a yes/no correctness prompt. Because the
judge is also the main reader, we keep it fixed when other models read \todo{re-grade a random subset with
gpt-6.1-sol and report agreement}. We use two answer prompts. The \emph{minimal} prompt tells the reader that
earlier sessions are no longer available, that only its memory of them is shown, and that it should reply as the
assistant using that memory. The \emph{tuned} prompt is the answer
prompt that Mem0 released for BEAM, which tells the reader to scan all memories, cross-check them, and use a fixed
sentence for abstention. On BEAM-1M it raises Hindsight by \prelim{8.6} points but our compiled notes by at most
\prelim{2.3} (\cref{sec:ablation}), so we report BEAM comparisons under the tuned prompt and name the prompt wherever it
changes.

% S4: baselines and budgets
\paragraph{Systems and budgets.}
Our main baseline is Hindsight~\citep{latimer2025hindsight}, which reports the highest BEAM scores among systems with
released code that we could run. We use its server (version 1.0.0) with gpt-6-luna for fact extraction; its recall
returns about 124K characters per question. Hindsight reports \prelim{73.9} on BEAM-1M with its own reader and judge;
under our protocol it scores \prelim{70.4} with the tuned prompt. Because \method{} makes read-time calls besides the
answer, we also run Hindsight\,+\,CC, which adds one call: a conflict check that lists retrieved memories that
disagree with each other and appends them to the context. \method{} still makes one more call than Hindsight\,+\,CC
(query planning, which the record-retrieval baseline also uses); \cref{sec:cost} reports the tokens each system
spends. We further report record retrieval, a re-implementation of LycheeMemory~\citep{li2026lychee} over the same
records that \method{} stores, with its default 8K-character context and, in \cref{sec:audit}, budgets up to 124K
characters; plain retrieval of raw dialogue chunks; and, on BEAM-100K, the full conversation. Embeddings come from
Qwen3-Embedding-8B, and hybrid search fuses dense and BM25 rankings by reciprocal rank. \Cref{tab:systems} lists the
calls and context length of each system.

% S5: statistics
\paragraph{Repetition and statistics.}
Every condition is run three times, with independent generations of every read-time LLM call; the store is built
once. We average the three scores of each question and compare systems on the same questions. Questions within a
conversation are correlated, so we resample whole conversations (5,000 bootstrap draws) and report 95\% percentile
intervals for differences between systems. On the 200 BEAM-100K development questions, the 95\% interval of a
single-generation score spans roughly $\pm$5 points, so we draw conclusions from three-run averages and mark every
single-run result.

\begin{table}[t]
\centering\small
\caption{Systems compared. ``Calls'' are LLM calls; ``ctx'' is the average context, in characters, that the reader
receives on BEAM-1M. $^\dagger$On BEAM-100K, the only tier where the conversation fits. \method{}'s context is the
notes plus the retrieved records. \todo{report tokens as well as characters; Hindsight write-time tokens.}}
\label{tab:systems}
\begin{tabular}{lllr}
\toprule
System & Write-time calls & Read-time calls & ctx \\
\midrule
Full context & none & answer & \prelim{675K}$^\dagger$ \\
Raw-chunk retrieval & none (embeddings) & answer & \prelim{7.4K} \\
Record retrieval (LycheeMemory) & 1 per segment & plan, answer & \prelim{7.6K} \\
Hindsight & fact extraction, entity summaries & answer & \prelim{124K} \\
Hindsight\,+\,CC & same & conflict check, answer & \prelim{124K} \\
\method{} & 1 per segment & plan, compile, answer & \prelim{64K} \\
\bottomrule
\end{tabular}
\end{table}
